\chapter{GİRİŞ}
\label{chap:giris}

\section{Araştırma Problemi ve Motivasyon}
\label{sec:giris-problem}

Artan talep, fosil kaynakların çevresel etkileri ve enerji arzının karbonsuzlaştırılması gereksinimi, elektrik enerjisi sistemlerinde önemli bir dönüşüme yol açmaktadır. Güneş ve rüzgâr gibi yenilenebilir kaynakların kurulu güç içindeki payı büyüdükçe elektrik üretimi, merkezi ve öngörülebilir bir yapıdan dağıtık ve zamana bağlı değişen bir yapıya doğru ilerlemektedir. Üretim yapısındaki dönüşüm temiz enerji üretimini artırsa da üretim ile tüketimin her an dengelenmesi zorunluluğunu ortadan kaldırmamaktadır. Aksine, ışınım ve rüzgâr koşullarına bağlı üretim değişimleri, güç dengesini korumak, gerilim ile frekansı desteklemek ve yenilenebilir enerjiyi etkin biçimde kullanabilmek için gereken esnek kaynaklara duyulan gereksinimi artırmaktadır \citep{Chen2009,Aneke2016,IEAPVPS2024}.

Enerji depolama sistemleri (EDS), enerjinin üretildiği anda tüketilmesi zorunluluğunu azaltarak yenilenebilir kaynakların şebekeye entegrasyonu için gereken esnekliği sağlayan başlıca yapılardan biridir. Bir EDS yalnızca batarya grubundan ibaret olmayıp batarya, güç dönüşüm sistemi, batarya ve enerji yönetim birimleri ile koruma, haberleşme ve ısıl yönetim alt sistemlerinin birlikte çalıştığı bütünleşik bir sistemdir \citep{Hannan2017,Plett2015}. Dolayısıyla sistem başarısını yalnızca hücrenin enerji yoğunluğu değil, batarya ile DA bara, yük ve şebeke arasındaki enerji akışını yöneten güç elektroniği arayüzünün verimi, güvenilirliği, hacmi ve maliyeti de belirlemektedir. Nitekim güncel maliyet çalışmalarında en büyük kalemi batarya paketi oluştursa da çift yönlü güç dönüştürücüleri ve yardımcı donanımları, sistem düzeyindeki maliyet ve kayıplarda önemli bir paya sahiptir \citep{Ramasamy2023,Lazard2024}. Söz konusu maliyet ve kayıp dağılımı, hem bataryanın güvenli çalışma bölgesini koruyan hem de aynı yarı iletken donanımı farklı enerji akış yönlerinde daha etkin kullanan dönüştürücü mimarilerini önemli kılmaktadır.

Depolama arayüzünün performansı ise yalnızca anma koşullarında değil, bağlı olduğu yenilenebilir üretim sisteminin gerçek çalışma profiliyle birlikte değerlendirilmelidir. Fotovoltaik (FV) üretim ışınım ve sıcaklığa bağlı olarak gün boyunca sürekli değiştiğinden, FV evirici anma gücüne ancak sınırlı sürelerle yaklaşmakta ve çalışma ömrünün kayda değer bir bölümünü kısmi yük bölgelerinde geçirmektedir. Bu nedenle anma noktasında verilen tepe verim değeri, yıllık enerji dönüşüm başarısını tek başına açıklamaya yetmemektedir. Avrupa ve CEC ağırlıklı verim tanımları ile bölgesel ışınım dağılımlarına dayalı ölçütler de bu sınırlılığı gidermek amacıyla farklı güç seviyelerindeki verimleri gerçek işletme koşullarındaki ağırlıklarına göre birlikte değerlendirmektedir \citep{Meinhardt2000,Sera2008,Sandia2015,Ongun2013}. Öte yandan kısmi yük bölgesi, verimin yanı sıra akım kalitesi açısından da kritiktir. Düşük güç seviyelerinde sabit sürme, anahtarlama ve manyetik eleman kayıplarının aktarılan güce oranı büyüdüğünden verim düşmekte, akım dalga biçimi de anahtarlama elemanlarının doğrusal olmayan davranışına, ölü zamana ve şebeke gerilimindeki bozulmalara karşı daha duyarlı hâle gelmektedir. Dolayısıyla düşük güç bölgesi, verim ile akım toplam harmonik bozulması (Total Harmonic Distortion, THD) arasındaki ödünleşimin en belirgin biçimde ortaya çıktığı çalışma alanıdır \citep{Kjaer2005,Chicco2009}.

Şebekeye bağlı dönüştürücülerde verim--THD ödünleşimini topoloji, modülasyon, filtre ve kontrol yapısı birlikte belirlemektedir. Topoloji düzeyinde, iki seviyeli eviricilere alternatif olarak önerilen çok seviyeli yapılar, çıkış gerilimini daha fazla seviyede oluşturarak yarı iletken gerilim zorlanmasını ve anahtarlama dalgalanmasını azaltabilmektedir. Çok seviyeli topolojilerden biri olan üç seviyeli T-tipi evirici, nötr noktaya bağlanan çift yönlü anahtar kolu sayesinde $+V_{\mathrm{dc}}/2$, $0$ ve $-V_{\mathrm{dc}}/2$ gerilim seviyelerini üretmektedir. Uygun eleman seçimi ve modülasyonla yüksek güç yoğunluğu ve yüksek verim potansiyeli sunan T-tipi yapı, ticari yarı iletken modüllerle de gerçekleştirilebilmektedir \citep{Wang2020}. Ayrıca T-tipi topolojinin şebekeye bağlı gerçek bir donanım üzerindeki uygulanabilirliği ve elektromanyetik uyumluluk gereksinimleri deneysel çalışmalarla gösterilmiştir \citep{Hizarci2021}.

T-tipi topolojinin anma gücündeki üstünlüğü ise düşük güçte aynı ölçüde korunmayabilir. Çıkış akımının küçük olduğu bölgelerde, T-tipi faz kolunda sıfır gerilim durumunu oluşturan nötr anahtarlar, akımın işaret değiştirmesine izin veren senkron bir alçaltıcı dönüştürücü (synchronous buck) davranışı gösterebilir. Böyle bir çalışmada endüktör akımı bir anahtarlama periyodu içinde sıfıra ulaştıktan sonra nötr anahtar iletimde tutulursa, akım ters yönde büyüyebilir ve şebekeye yararlı güç taşımayan bir dolaşım akımı oluşabilir. Ortaya çıkan dolaşım akımı, yarı iletken ve pasif elemanların RMS akımı ile iletim kayıplarını artırabilir. Benzer bir sorunla karşılaşılan senkron alçaltıcı dönüştürücülerde, süreksiz iletim ve diyot benzetimi (diode emulation) yaklaşımlarının hafif yükte ters akımı engelleyerek kayıpları azaltabildiği bilinmektedir \citep{Erickson2001,Mohan2003,Chen2010}. Ancak ters akımı engelleme ilkesini şebekeye bağlı üç seviyeli T-tipi eviricinin sıfır durumuna uygulamak, yalnızca bir anahtarı kesime götürmekten ibaret değildir. Sıfır durumunda yapılan değişiklik akım yolunu, çıkış gerilimini, modülasyon bütünlüğünü ve harmonik içeriği birlikte etkilediğinden, problem T-tipi topolojiye özgüdür.

Tezin ilk odak noktası, T-tipi sıfır durumuna özgü ters akım problemi için geliştirilen yeni bir anahtarlama yöntemidir. Yöntemin ilkesi, düşük güç bölgesinde dolaşımı sınırlandırmak amacıyla nötr anahtarın belirli koşullarda susturulmasına dayanmaktadır. Susturma, dolaşım akımını ve bu akımın yol açtığı kayıpları azaltmayı hedeflerken T-tipi eviricinin ürettiği gerilim darbelerini de değiştirmektedir. Gerilim darbelerindeki söz konusu değişim nedeniyle, verimdeki olası iyileşmenin karşılığında özellikle üçüncü harmonik bileşeninde artış ortaya çıkabilmektedir. Üstelik ölü zaman, anahtar geçişleri ve parazitik kapasiteler gibi gerçek devre etkileri de akım harmoniklerini çalışma noktasına bağlı biçimde değiştirebilmektedir \citep{Zhang2014Deadtime}. Bu nedenle önerilen anahtarlama yöntemi, verim kazancının yanında aynı çalışma noktalarındaki THD ve şebeke akımı kalitesiyle birlikte değerlendirilmelidir.

Anahtarlama değişikliğinin oluşturduğu üçüncü harmonik bileşenini azaltmak için tez kapsamında Sin\textsuperscript{3} tabanlı bir ileri besleme kompanzasyonu (feedforward compensation) kullanılmaktadır. Kompanzasyonun temelini oluşturan $\sin^3(\theta)=[3\sin(\theta)-\sin(3\theta)]/4$ özdeşliği, ölçülen şebeke geriliminin faz bilgisinden temel ve üçüncü harmonikle ilişkili anlık bir şekillendirme terimi üretilebilmesini sağlamaktadır. Şebeke gerilimi ileri besleme yolları akım bozulmasını azaltabilse de, dijital gecikme, LCL filtre ve şebeke empedansı üzerinden sistem dinamiğiyle etkileşebildiklerinden kararlılık açısından incelenmeleri gerekmektedir \citep{Lu2016,Wang2024}.

Tezin ikinci odak noktası, aynı T-tipi donanımın evirici işlevinin yanında EDS'nin öteki enerji akış kiplerinde de yeniden kullanılabilmesidir. Geleneksel bir depolama sisteminde şebeke eviricisi, şarj doğrultucusu veya güç faktörü düzeltme (Power Factor Correction, PFC) katı ve hücre dengeleme dönüştürücüleri ayrı yarı iletken gruplarıyla gerçekleştirilebilir; ancak her ilave kat güç yoğunluğu, sürücü, soğutma, koruma ve maliyet üzerinde yeni bir yük oluşturmaktadır. Tezde önerilen GT-tipi mimari ise ayrı katlar yerine T-tipi dönüştürücünün anahtarlarını çalışma kipine göre farklı görevlerde kullanarak evirici, Vienna PFC ve aktif batarya yönetimi işlevlerini ortak bir donanım üzerinde birleştirmeyi amaçlamaktadır. Bu amaçla T-tipi dönüştürücünün yapısal olanaklarından iki ölçekte yararlanılmaktadır: faz kolu ölçeğinde nötr kolun düşük güçteki görevi değiştirilmekte, sistem ölçeğinde ise dönüştürücü anahtarları şarj--deşarj mimarisi içinde yeniden görevlendirilmektedir.
\section{Tezin Savı, Amacı ve Araştırma Soruları}
\label{sec:giris-amac}

Tezin ana savı iki önermeden oluşmaktadır: \emph{T-tipi dönüştürücünün yapısal ve anahtarlama özellikleri, enerji depolama sistemlerinde şarj--deşarj işlevlerini yeniden yapılandırmak ve düşük güçte verim--harmonik ödünleşimini iyileştirmek için birlikte kullanılabilir. Sin\textsuperscript{3} tabanlı ileri besleme kompanzasyonu, düşük güç verim--harmonik hedefini desteklemektedir.} Literatürde T-tipi eviricinin yüksek güç yoğunluğu, şebeke kontrolü, hafif yük dönüştürücü kipleri ve ileri besleme yolları ayrı ayrı incelenmiştir \citep{Blaabjerg2006,Wang2020,Chen2010,Lu2016}. Ancak T-tipi faz kolundaki düşük güç ters akım mekanizmasının, bu mekanizmanın verim--THD sonucunun, Sin\textsuperscript{3} tabanlı kompanzasyonun ve aynı donanımın EDS işlevleri için yeniden yapılandırılmasının ortak bir mimari ve doğrulama zinciri içinde birlikte incelenmesi, tezin ele aldığı araştırma boşluğunu oluşturmaktadır.

Tezin temel amacı, enerji depolama sistemleri için T-tipi dönüştürücü tabanlı yeniden yapılandırılabilir bir şarj--deşarj mimarisi geliştirmektir. Temel amacın yanında, aynı topolojinin şebekeye bağlı düşük güç çalışmasında görülen verim ile harmonik bozulma sorunlarını birlikte iyileştirecek yöntemlerin ortaya konması da amaçlanmaktadır. Bu amaçlar doğrultusunda belirlenen araştırma soruları aşağıda sıralanmıştır:

\begin{enumerate}[label=AS\arabic*.]
  \item T-tipi eviricinin düşük güç çalışma bölgesinde ters akım dolaşımı hangi anahtarlama durumları ve akım yolları üzerinden oluşmakta, oluşan dolaşım akımı kayıpları ve akım harmoniklerini nasıl etkilemektedir?
  \item Nötr anahtarın koşula bağlı olarak susturulması, klasik anahtarlamaya göre düşük güç verimini hangi çalışma aralığında ve hangi harmonik maliyetle değiştirmektedir?
  \item Sin\textsuperscript{3} tabanlı ileri besleme kompanzasyonu, önerilen anahtarlamanın oluşturduğu üçüncü harmonik bileşenini ne ölçüde azaltmakta ve kapalı çevrim akım denetiminin kararlılık marjlarıyla nasıl etkileşmektedir?
  \item T-tipi dönüştürücünün yarı iletkenleri, EDS'nin şarj--deşarj ve batarya bazlı şarj denetimi kiplerinde güvenli biçimde yeniden görevlendirilerek ayrı güç dönüşüm katlarının işlevleri ortak bir yapıda birleştirilebilir mi?
\end{enumerate}

\section{Özgün Katkılar}
\label{sec:giris-katkilar}

Tezin özgün katkıları, ortak tabanını T-tipi topolojinin oluşturduğu iki ölçekte toplanmaktadır: ilk iki katkı T-tipi faz kolu ölçeğinde, üçüncüsü ise sistem ölçeğindedir.

\begin{enumerate}
  \item \textbf{Düşük güç anahtarlama yöntemi:} T-tipi eviricinin sıfır gerilim durumunda ortaya çıkan ters akım dolaşımını hedefleyen bir nötr anahtar susturma yöntemi geliştirilmiştir. Yöntem, çalışma kiplerinin ve akım yollarının analiziyle birlikte benzetim ve prototip deney kayıtları üzerinden değerlendirilmektedir. Düşük güç deneyinde susturma verimi 1,7 puan artırmış, ancak bedel olarak akım THD'sini yükseltmiştir.
  \item \textbf{Sin\textsuperscript{3} tabanlı ileri besleme kompanzasyonu:} Yeni anahtarlama yönteminin üçüncü harmonik üzerindeki etkisini azaltmak üzere, şebeke işaretinden anlık olarak üretilen bir kompanzasyon terimi önerilmiştir. Akım çevriminin modülasyon girişine eklenen bir ileri besleme bileşeni olarak modellenen terim, çevrimin karakteristik denklemini değiştirmemektedir. Susturma ile birlikte kullanıldığında terim, deneyde akım THD'sini klasik anahtarlama düzeyine indirmiş ve verimi korumuştur.
  \item \textbf{GT-tipi yeniden yapılandırılabilir mimari:} Evirici, Vienna PFC ve aktif batarya yönetimi işlevlerinde aynı yarı iletkenlerin farklı görevlerle kullanılmasını sağlayan T-tipi tabanlı bir EDS mimarisi önerilmiştir. Önerilen yapı, klasik ayrık mimariyle çalışma kipleri, anahtarlama durumları, enerji akış yolları ve güvenlik kısıtları üzerinden karşılaştırılmaktadır. Deney düzeneği GT-tipi olarak kurulmuş olup donanımın evirici kipi düşük güç deneyleriyle, şarj yönündeki işleyişi ise benzetim kayıtlarıyla gösterilmiştir. Mimari ayrıca, ayrı güç katlarına göre eleman sayısını azaltma olanağı sunmaktadır.
\end{enumerate}

Tez başlığındaki ``enerji depolama sistemleri için geliştirilmiş T-tipi dönüştürücü mimarisi'' GT-tipi mimariyi (katkı~3), ``düşük güç performansını iyileştiren anahtarlama'' nötr anahtar susturmayı (katkı~1), ``Sin\textsuperscript{3} tabanlı kompanzasyon'' ise ileri besleme kompanzasyonunu (katkı~2) adlandırmaktadır. Başlıktaki sıralamadan farklı olarak bölümlerin sırası fiziksel akışı izlemektedir: önce faz kolu ölçeğindeki mekanizma ve iki yöntem, ardından aynı donanımın sistem ölçeğinde yeniden görevlendirilmesi incelenmektedir.

Katkı~1 ve katkı~2'de aynı S1--S4 benzetim modelleri kullanılmıştır. Modellerin doğrulanması Alt Bölüm~\ref{sec:model-dogrulama-duzeltme} içinde verilmiş, bütün parametreler kaynak ve tür bilgisiyle birlikte tek tabloda toplanmıştır (Tablo~\ref{tab:sistem-parametreleri}). Modeller deneydeki aktif ve reaktif güce göre kalibre edilmiş, öteki büyüklükler ise deneyle bağımsız olarak karşılaştırılmıştır (Alt Bölüm~\ref{sec:kalibrasyon}).

\section{Tezin Kapsamı}
\label{sec:giris-kapsam}

Çalışma, şebekeye bağlı üç seviyeli T-tipi dönüştürücünün düşük güç davranışını ve T-tipi topolojiden türetilen GT-tipi enerji depolama mimarisini kapsamaktadır. Kurulan prototip, nötr hattı DA bara orta noktasına bağlı üç faz, dört telli bir yapıdadır. Benzetim modelleri ve deney değerlendirmesi ise nötr dönüşü yine DA bara orta noktasına bağlanan tek bir faz kolu üzerinde yürütülmüştür. Düşük güç incelemesinde anahtarlama durumları, ters akım, dönüştürücü verimi, şebeke akımı THD'si ve üçüncü harmonik temel performans ölçütleri olarak ele alınırken GT-tipi mimari, çalışma kipleri, enerji akış yolları ve güvenli kip geçiş koşulları düzeyinde incelenmektedir.

Doğrulama, dört senaryonun benzetim modelleri ve düşük güç noktasında yapılan deneylerle gerçekleştirilmiştir. Ele alınan senaryolar klasik tamamlayıcı anahtarlama (S1), nötr anahtar susturma (S2), Sin\textsuperscript{3} kompanzasyonlu susturma (S3) ve akım geri beslemeli kapalı çevrim benzetimidir (S4). S1--S3 deneyleri, şebeke gerilimi varyak üzerinden yaklaşık \SI{120}{\volt} RMS değerine ayarlanarak güvenli laboratuvar koşulunda ve yaklaşık \SI{2.3}{\ampere} faz akımında yürütülmüştür. Anma akımı koşulu ise \SI{25}{\ampere} hedefli S4 benzetimiyle incelenmiş olup donanımın anma gücü değerleri tasarım hedefi olarak verilmektedir. Sonuçlar, tezde tanımlanan DA bara gerilimi, şebeke koşulları, LCL filtre, anahtarlama frekansı, denetleyici ayarları ve ölçüm düzeninde elde edilmiştir.

\section{Tezin Düzeni}
\label{sec:giris-duzen}

Tez sekiz bölümden oluşmakta ve genelden özele, faz kolundan sisteme doğru ilerlemektedir. Problemden katkılara uzanan bu akış ile her katkının incelendiği bölüm Şekil~\ref{fig:tez-akisi} üzerinde verilmiştir.

\begin{figure}[H]
  \centering
  \resizebox{\textwidth}{!}{%
  \begin{tikzpicture}[
    node distance=6mm and 7mm,
    kutu/.style={draw, rounded corners=1.5pt, align=center, font=\small, minimum height=12mm, text width=32mm, fill=blue!5, execute at begin node=\nohyph},
    sorun/.style={kutu, fill=red!8},
    katki/.style={kutu, fill=orange!15, very thick},
    bedel/.style={kutu, fill=gray!10},
    kanit/.style={kutu, fill=green!8},
    akis/.style={-{Latex[length=2mm]}, thick},
    ince/.style={-{Latex[length=1.6mm]}, dashed},
    grup/.style={draw, dashed, rounded corners=3pt, inner sep=3.5mm}
  ]
    \def\nohyph{\hyphenpenalty=10000\relax}
    \node[kutu] (eds) {Enerji depolama sistemi:\\güç dönüşüm arayüzü};
    \node[kutu, right=of eds] (tt) {Üç seviyeli T-tipi\\dönüştürücü\\(Bölüm~\ref{chap:literatur}--\ref{chap:sistem})};
    \node[sorun, right=of tt] (ters) {Düşük güçte nötr kol\\üzerinden ters akım\\dolaşımı};
    \node[kutu, right=of ters] (odun) {Verim--THD\\ödünleşimi};
    \draw[akis] (eds) -- (tt);
    \draw[akis] (tt) -- (ters);
    \draw[akis] (ters) -- (odun);

    \node[katki, below=14mm of ters] (k1) {\textbf{Katkı 1}\\Nötr anahtar susturma\\(Bölüm~\ref{chap:anahtarlama})};
    \node[bedel, right=of k1] (dcm) {Bedel: süreksiz iletim,\\üçüncü harmonik};
    \node[katki, below=of dcm] (k2) {\textbf{Katkı 2}\\{Sin\textsuperscript{3}} ileri besleme\\kompanzasyonu (Bölüm~\ref{chap:sin3})};
    \node[kanit, left=of k2] (kc) {Kapalı çevrim S4\\ve kararlılık analizi\\(Alt Bölüm~\ref{sec:s4-kapali-cevrim}--\ref{sec:kararlilik-bode})};
    \node[kanit, below=of kc] (dny) {Deney kayıtları ve\\benzetim modeli\\(Bölüm~\ref{chap:deney})};
    \draw[akis] (ters) -- (k1);
    \draw[akis] (k1) -- (dcm);
    \draw[akis] (dcm) -- (k2);
    \draw[ince] (k2) -- (kc);
    \draw[ince] (k1.west) -- ++(-5mm,0) |- (dny.west);
    \draw[ince] (k2.south) |- (dny.east);
    \node[grup, fit=(k1)(dcm)(k2)(kc)(dny), label={[font=\small\itshape]above right:Faz kolu ölçeği}] (faz) {};

    \node[katki, below=14mm of tt, xshift=-20mm] (k3) {\textbf{Katkı 3}\\GT-tipi yeniden\\yapılandırılabilir mimari\\(Bölüm~\ref{chap:gttipi})};
    \node[bedel, below=of k3] (isl) {Evirici, Vienna PFC ve\\aktif batarya yönetimi\\aynı yarı iletkenlerle};
    \draw[akis] (tt.south) -- ++(0,-5mm) -| (k3.north);
    \draw[ince] (k3) -- (isl);
    \node[grup, fit=(k3)(isl), label={[font=\small\itshape]above left:Sistem ölçeği}] (sis) {};

    \node[kutu, anchor=north west, text width=96mm] (taban) at ([yshift=-8mm]sis.south west |- faz.south) {Ortak taban: sistem modeli, tek parametre tablosu ve benzetim modelinin doğrulanması (Bölüm~\ref{chap:sistem})};
    \node[kutu, right=of taban, text width=44mm] (son) {Araştırma sorularının yanıtları,\\gelecek çalışmalar (Bölüm~\ref{chap:sonuc})};
    \draw[ince] (taban.north -| sis.south) -- (sis.south);
    \draw[ince] (taban.north -| dny.south) -- (faz.south -| dny.south);
    \draw[akis] (son.north |- faz.south) -- (son.north);
  \end{tikzpicture}%
  }
  \caption[Tezin problem--yöntem--katkı akışı]{Tezin problem--yöntem--katkı akışı.}
  \label{fig:tez-akisi}
\end{figure} 


Birinci bölümde araştırma problemi, amaç, araştırma soruları, katkılar ve kapsam tanımlanmıştır. Ardından ikinci bölüm, EDS güç dönüşüm mimarilerine, T-tipi eviricilere, akım kontrolüne ve ileri besleme yaklaşımlarına ilişkin literatürü inceleyerek araştırma boşluğunu belirlemektedir. Üçüncü bölüm ise sistem modelini, parametreleri ve benzetim modelinin doğrulanmasını sunarak sonraki bölümlerin ortak tabanını oluşturmaktadır.

Faz kolu ölçeğindeki incelemeye dördüncü bölümle geçilmektedir. Dördüncü bölüm düşük güç ters akım mekanizmasını ve nötr anahtar susturma yöntemini incelemektedir (katkı~1). Beşinci bölüm, susturmanın bedeli olarak ortaya çıkan üçüncü harmoniğe yönelik Sin\textsuperscript{3} tabanlı ileri besleme kompanzasyonunu, kapalı çevrim modelini ve kararlılık incelemesini sunmaktadır (katkı~2). Altıncı bölüm, iki yöntemi prototipin deney kayıtlarıyla karşılaştırmakta ve benzetim modelini deney özetine kalibre etmektedir. Yedinci bölüm ise sistem ölçeğine geçerek GT-tipi yeniden yapılandırılabilir EDS mimarisini ve çalışma kiplerini sunmaktadır (katkı~3). Son olarak sekizinci bölümde araştırma soruları yanıtlanmakta, özgün katkılar değerlendirilmekte ve gelecek çalışmalar önerilmektedir.
